\chapter{Background and related work}\label{ch:background}

\section{Linguistic ingroup bias on Wikipedia}\label{sec:wikipedia_bias}

\citet{oeberst2020wikipedia} establish the load-bearing prior for this
thesis. Across thirty-five intergroup conflicts and three studies, they
show that Wikipedia articles in the language of each conflicting party
portray the in-group as preferred, more powerful, less immoral, and less
responsible, using a four-dimensional content-coding rubric applied to
article lead sections. The strength of the bias is predicted by two
structural variables: the recency of the conflict, and the proportion of
in-group editors among the article's top contributors. A null result on
translation, rater nationality, and scale-order licenses subsequent
embedding-based work that does not require native-speaker raters.
Because the biased unit is the human-authored article text --- the
same corpus that contemporary multilingual LLMs consume in
pre-training --- their result fixes the source-side prior for this
thesis: a model that reproduces the distribution of its training text
on contested events should reproduce the per-language framing
divergence along with it. Whether it does, and in what form, is the
question the following chapters operationalise;
\S\ref{sec:mechanisms} returns to this as the source-side mechanism.

\citet{samoilenko2017timelines} document a related effect at corpus
scale. They build a language-independent pipeline that extracts every
four-digit-year mention from ``History of [country]'' articles across
thirty Wikipedia editions for all 193 UN member states. Using a
Dirichlet--multinomial null model, Jensen--Shannon divergence, and
hierarchical clustering, they find that national histories cluster by
geopolitical bloc, that recency is over-represented in every edition (a
uniform ``presentism''), and that the post-Soviet bloc surfaces as a
coherent cluster of editions sharing temporal focal points. The
(country $\times$ language-edition) comparison and clustering frame is
methodologically analogous to the (question $\times$ language) anchor
matrices used in this thesis. Their explicit non-goal --- language-specific
narrative and sentiment analysis --- is the gap that the embedding-based
contested-event method fills here.

\citet{yasseri2026grokipedia} extend the cross-corpus comparison from
human-edited Wikipedia to LLM-generated encyclopedic text. Comparing
17{,}790 matched Grokipedia--Wikipedia pairs from the most-edited English
Wikipedia titles, they find that Grokipedia articles are roughly 22\%
longer with 43\% fewer references per thousand words, with a distinctly
bimodal similarity distribution: roughly 34\% of pairs are aligned and
66\% diverge. The divergent subset shows a systematic rightward shift
in the political bias of cited news media, concentrated in articles on
history, religion, and the arts. They name the resulting phenomenon the
\emph{epistemic opacity paradox}: Grokipedia \emph{appears} neutral
because it lacks visible disagreement, but its bias is hidden in opaque
model behaviour rather than being visible and contestable as in
Wikipedia's edit history. The bimodal pattern --- bias concentrated in
narrative-laden content --- anticipates the finding in
Chapter~\ref{ch:experiments} that cross-lingual divergence in LLM
responses is content-dependent rather than uniform across questions.

A complementary structural lens comes from
\citet{yu2025demographic}, who decompose the coverage of over one
million biographies across the twelve largest editions (2010--2020)
into \emph{breadth}, \emph{depth}, and \emph{global consensus}, and
isolate editorial-system effects from upstream fame asymmetry with
an audit on Olympic medal winners. The decomposition is portable to
characterising the multilingual anchor corpus, with the caveat that
their cohort does not include Ukrainian.

Earlier work on cross-lingual Wikipedia divergence sets the stage.
\citet{hecht2010tower} document the ``tower of Babel'' effect: language
editions of Wikipedia overlap less than naive intuition suggests, with
significant article-level and content-level differences across editions
even on shared topics. \citet{massa2013manypedia} introduce
\emph{Manypedia}, the first systematic comparison of language-specific
points of view on the same topic across Wikipedia communities. Both
papers establish that \emph{which language one reads Wikipedia in}
materially shapes the narrative one receives --- a precondition for any
claim about language-conditional LLM behaviour.

The information-ecosystem framing is rounded out by
\citet{pfister2011networked} on participatory expertise and
\citet{urman2022shadows} on the linguistic siloing of the
multilingual far-right web on Telegram: contemporary online
information ecosystems are linguistically partitioned by design.
Neither offers methods that this thesis inherits.

\section{Embedding-space bias measurement}\label{sec:embedding_bias}

The use of cosine geometry to detect bias in dense vector representations
begins with \citet{bolukbasi2016man}, who identify a single ``gender
direction'' in word2vec via principal-components analysis on paired
\emph{she--he} and \emph{woman--man} difference vectors and demonstrate
that stereotyped analogies (e.g.\ \emph{programmer $\leftrightarrow$
man, homemaker $\leftrightarrow$ woman}) drop from 19\% to 6\% under
their neutralize-and-equalize debiasing procedure with no loss on
standard analogy benchmarks. \citet{caliskan2017semantics} formalise the
cosine-association methodology as the Word Embedding Association Test
(WEAT) --- measuring the differential association of two target word
sets with two attribute word sets via cosine similarity, with a
permutation-test null --- and show that pre-trained GloVe embeddings on
Common Crawl encode every IAT-documented human bias they tested, from
flowers/insects ($d = 1.50$) to gender--career ($d = 1.81$) and
Bertrand--Mullainathan r\'esum\'e-name discrimination ($d = 1.50$,
$p < 10^{-4}$). Their WEFAT extension predicts real-world quantities:
the gender composition of fifty occupations correlates with embedding
geometry at $\rho = 0.90$ against 2015 BLS data, with effects robust
across GloVe and word2vec at $\rho = 0.88$.

{\sloppy
\citet{may2019} extend WEAT from word level to sentence level (SEAT),
demonstrating that contextualised sentence encoders reproduce many of
the same effects when target and attribute words are embedded in
templated sentences. \citet{kurita2019measuring} push further into the
contextualised regime, proposing a template-based log-probability bias
score for masked language models: insert a target word and an attribute
word into a templated sentence, query the masked-LM head for both, and
compute $\log(p_\text{target} / p_\text{prior})$ where the prior is the
model's probability of the target with the attribute also masked.
Critically, they show that cosine-based WEAT applied to BERT
\emph{fails} to detect any of the five tested biases at $p < 0.01$,
while their log-probability score recovers statistically significant
bias on four of five categories and aligns better with human IAT
measurements. The lineage is now clear: word-level cosine
\citep{bolukbasi2016man, caliskan2017semantics} $\rightarrow$
sentence-level cosine \citep{may2019} $\rightarrow$ masked-token log
probability in templated sentences \citep{kurita2019measuring}.\par}

The methodology used in this thesis is the natural next step along the
same lineage: measure bias at the \emph{response} level --- full
multi-sentence LLM outputs --- by computing cosine similarity to
language-anchored Wikipedia lead sections embedded in the same space.
The shift from stereotype-direction measurement (the WEAT family) to
\emph{ingroup pull} measurement (the cosine of an LLM response to a
language-specific anchor) is conceptually a contextual-WEAT effect with
response-cells as targets and language-anchored Wikipedia leads as
attributes; Caliskan's permutation-test null transfers directly. The
bias-subspace projection idea introduced by \citet{bolukbasi2016man}
also transfers: it motivates a ``Russian-versus-English framing axis''
constructed via PCA on paired EN/RU/UK Wikipedia-anchor difference
vectors, projection onto which provides a one-dimensional summary
complementary to the cosine-to-anchors instrument
(\S\ref{sec:debiasing}).

\citet{azizov2024safari} provide an independent validation of Wikipedia
descriptions as effective cross-lingual anchors for political-bias and
factuality detection across ten or more languages. Their finding that
nationality bias is prevalent when local Wikipedia context is missing
motivates the careful per-language anchor construction used in
Chapter~\ref{ch:method}.

\section{Multilingual LLMs as knowledge curators}\label{sec:llm_curators}

The empirical literature on bias in conversational LLMs begins with
\citet{hartmann2023chatgpt}, who prompt ChatGPT with 630 political
statements drawn from Germany's Wahl-O-Mat, the Netherlands' StemWijzer,
and the political-compass test, scoring its forced
agree/neutral/disagree replies through the voter-advice-application
backends and finding a consistent pro-environmental, left-libertarian
orientation across all three pre-registered studies, robust to five
prompt manipulations including English--Spanish translation.
\citet{rozado2024} extends the analysis across twenty-four LLMs and
reports that the same left-of-centre orientation is shared across nearly
all conversational models tested. Notably, Rozado reports that
\emph{base} models cluster near the political centre indistinguishably
from a random-answer baseline, while conversational variants of the
same models exhibit the left-of-centre tilt, suggesting that the
political orientation emerges via post-pretraining alignment
(SFT/RLHF) rather than being induced by the pretraining corpus.
Rozado himself flags this as ``suggestive but ultimately inconclusive''
because the high invalid-response rate of base models (mean 42\%) and
moderate human--automated stance-agreement on those responses
($\kappa = 0.41$) leave the base-vs-conversational comparison
underpowered. \citet{dang2024rlhf}
supply the missing causal mechanism: preference optimization (DPO/RLOO)
applied with \emph{English-only} preference data demonstrably improves
model behaviour in twenty-two other languages on Aya-23-8B, showing
that RLHF signal leaks across the language barrier. Together, Rozado
and Dang predict that an English-data-aligned model will surface its
alignment-data tilt in Russian and Ukrainian outputs too --- a
mechanistic candidate for the EN-ingroup pull reported in
Chapter~\ref{ch:experiments}, complementary to \citet{zhao2025makieval}'s
English Activation effect. The methodology of \citet{hartmann2023chatgpt}
and \citet{rozado2024} --- forced-choice voter-advice questionnaires
projected onto a low-dimensional political space --- is sufficient for
the issue-position question but not for the contested-event question
pursued here.

\citet{durmus2023globalopinionqa} provide the precedent that motivates
this thesis's design directly. Their \emph{GlobalOpinionQA} dataset
combines 2{,}556 questions from the Pew Global Attitudes Survey and the
World Values Survey, and their per-country alignment metric is one minus
the Jensen--Shannon divergence between the model's option distribution
and the empirical option distribution from each country's respondents.
Across three experimental conditions --- default prompting,
cross-national prompting (``how would someone from $c$ answer''), and
\emph{linguistic prompting} (translating the question into Russian,
Chinese, or Turkish) --- they find three results. The model defaults to
opinions most similar to those of US, Canadian, Australian, and Western
European respondents; cross-national prompting shifts responses toward
the named country but often via shallow stereotype rather than nuanced
understanding; and crucially, \emph{linguistic prompting does not make
responses more similar to the speakers of the target language}. This
last result is the published precedent for the central empirical claim
of this thesis --- that prompting an LLM in Russian or Ukrainian does
not reliably elicit a Russian or Ukrainian framing of contested events
--- and it is generalised here from survey opinions to historical-event
narrative.

\citet{smirnov2026language} provides the closest published parallel
to the central experimental design of this thesis. On a single
Ukrainian civil-society document, Smirnov shows that semantically
equivalent prompts in Russian and Ukrainian elicit systematically
different ideological framings from a single LLM: the
Russian-language output reproduces vocabulary characteristic of
Russian state discourse (civil-society actors framed as
``illegitimate elites undermining democratic mandates''), while the
Ukrainian-language output adopts vocabulary characteristic of
Western liberal-democratic political science (the same actors framed
as ``legitimate stakeholders''). The model and source document are
held fixed; the prompt language alone flips the framing.
Smirnov's design is qualitative-discourse-analytic, single-document,
single-model (likely), and lacks an English control or an
imaginary-conflict probe; the present thesis is best read as the
embedding-quantified, English-controlled scaling of Smirnov's
existence proof across many events and providers.

\citet{buyl2024llm} provide the closest published comparable for the
cross-provider design used here. They prompt nineteen LLMs in six UN
languages to morally rate roughly four thousand Pantheon-listed political
figures via a two-stage protocol (open description, then forced
five-point Likert), aggregating responses into 61-dimensional
Manifesto-tag vectors and projecting to ideological axes via PCA. Their
headline finding is that the ideological position of an LLM depends
jointly on the language of the prompt and on the geopolitical region of
the model's creator, with significant \emph{within-bloc} spread (Gemini
progressive vs.\ Grok sovereigntist; Qwen international vs.\ ERNIE
domestic). They argue normatively, drawing on Foucault, Gramsci, and
Mouffe, that ``neutrality'' is itself ideologically defined, and that
regulatory policy should target ideological transparency and prevent
monoculture rather than enforce a single neutral position. The
methodological blind-spot they report --- silently filtering ``invalid''
responses --- is precisely the refusal signal that this thesis treats
as primary evidence (Chapter~\ref{ch:experiments}).
\citet{pacheco2025echoes} narrows the cross-provider comparison to
ChatGPT versus DeepSeek and finds that ``sovereign models'' reflect the
national power structures and narratives of their origin and narrow the
range of political thought in their respective regions, consistent with
the cross-provider divergence reported in
Chapter~\ref{ch:experiments}.

The cultural-bias literature is adjacent. \citet{ramezani2023moral}
probe English pre-trained language models with English-language
prompts about cross-cultural moral norms and document substantial
variation in LLMs' representations: the inferred norms align
substantially better with Western or Rich-West cultures than with
non-Western or non-rich cultures, even though the probe is held in
a single language throughout.
\citet{naous2024beer} introduce CAMeL, a probe for cultural bias in LLMs
in Arabic versus English, and report that even models with strong
Arabic capabilities default to Western-cultural framings on
culturally-contested items. \citet{atari2023humans} make the WEIRD
framing explicit: most LLM evaluations use Western, Educated,
Industrialised, Rich, Democratic populations as their reference,
leaving the cultural representation of the rest of the world
under-specified. \citet{zhao2025makieval} introduce MAKIEval, a
Wikidata-based framework for cultural-awareness evaluation, and report
an \emph{English Activation} effect --- English prompts often activate
a model's cultural knowledge better than the native language of that
culture. This is a direct mechanistic candidate for the EN-ingroup pull
reported in Chapter~\ref{ch:experiments}: prompting in English may
simply route through a richer, more densely cross-referenced
cultural-knowledge representation.

The cross-lingual factual-consistency literature provides the most
serious methodological challenge. \citet{qi2023crosslingual} introduce
\emph{RankC}, a softmax-weighted MAP-at-$K$ metric that quantifies
cross-lingual factual consistency \emph{independently of probing
accuracy}, and apply it to the \emph{BMLAMA} benchmark across XLM-R,
mT5, and the BLOOM series. Three findings here bear directly on this
thesis. First, average consistency across models is low (25--32\%) and
barely improves with scale (+2 percentage points from 560M to 3B).
Second, \emph{subword-vocabulary overlap} is the strongest predictor of
cross-lingual factual consistency --- far more than genetic, syntactic,
or geographic typology --- suggesting that factual knowledge percolates
between languages shallowly via shared tokens rather than via
language-agnostic representations. Third, in the BLOOM tokenizer the
Russian--Ukrainian pair has subword overlap of 0.76 and genetic
similarity of 0.8, placing it firmly among the high-CLC pairs. The
implication is that the finding in Chapter~\ref{ch:experiments} that
the Russian and Ukrainian Wikipedia anchors and both languages'
responses fuse into a single joint cluster is \emph{predicted} by
tokenizer overlap alone, and
the discussion in Chapter~\ref{ch:discussion} must explicitly
disentangle tokenizer leakage from genuine narrative fusion. The
proposed tokenizer-overlap control is reported in
Chapter~\ref{ch:experiments}: tokenizer leakage is ruled out as
the dominant explanation for the Ukrainian--Russian cluster fusion.

Two follow-up papers refine the tokenizer story specifically for the
Ukrainian side. \citet{maksymenko2025tokenization} measure tokenizer
fertility on Ukrainian across Llama, Gemma, Mistral, and the GPT
family, showing that the large Cyrillic vocabulary in these
tokenizers is dominated by Russian-frequent subwords that fail to
encode Ukrainian efficiently --- Gemma has more Cyrillic tokens than
Llama 3.1 yet produces \emph{worse} Ukrainian fertility because those
tokens are Russian-shaped. \citet{kiulian2024bilingual} report that
GPT-4 spends roughly three times more tokens on Ukrainian than on
English, and that this fragmentation mechanically drives
code-switching in vanilla Mistral on Ukrainian (code-switch rate
0.515), which targeted vocabulary expansion fixes (0.001). Neither
is directly portable (the pipeline is API-bound and cannot retrain
tokenizers), but together
with \citet{qi2023crosslingual} they tighten the chapter-five caveat:
high script-level overlap between Russian and Ukrainian masks
sub-token divergence that may itself contribute to the embedding-space
fusion pattern even after the tokenizer-overlap control rules out
the simple overlap-leakage story.

{\sloppy
\citet{urman2025silence} provide the closest published comparable for
refusal-as-evidence. Auditing ChatGPT (GPT-3.5/4/4o), Google
Bard/Gemini, and Bing Chat on Russia-related political prompts in
Russian, Ukrainian, and English, they find that Bard refuses 90\% of
Putin-related queries in Russian versus 19\% in English --- a refusal
pattern that closely mirrors a leaked Yandex/Roskomnadzor block list.
The methodological recipe --- manual two-annotator coding with Cohen's
$\kappa$ of 0.86--0.93, Fisher's exact tests with Benjamini--Hochberg
correction, and query-set seeding from a Russian state block list ---
is the gold-standard guardrail-audit recipe and the closest precedent
for the YandexGPT refusal sweep in Chapter~\ref{ch:experiments}. They
predict that Western chatbots will inherit Russian-state censorship
templates via guardrails; this thesis closes their Russian-provider gap
by showing that the Russian LLM enforces the templates uniformly (270
of 270 prompts) with an embedding-space signature visible via EVoC
clustering. The reframing of YandexGPT as the limit case of a
structural inheritance pattern --- rather than a categorically distinct
refusal regime --- is left to the discussion in
Chapter~\ref{ch:discussion}.\par}

\citet{makhortykh2024warriors} provide the answered-but-wrong companion
to \citet{urman2025silence}'s refusal study. From the same lab,
auditing Gemini, Bing/Copilot, and Perplexity on twenty-eight Kremlin
disinformation narratives in a multilingual longitudinal design, they
find 2024 inaccuracy rates spanning 10\% (Perplexity) to 48\% (Gemini),
with the critical observation that \emph{Gemini improves only in
English while deteriorating in low-resource languages including
Russian and Ukrainian}. Together \citet{urman2025silence} and
\citet{makhortykh2024warriors} bracket the LLM cross-lingual failure
modes: refusal on the silence side, narrative reproduction on the
content side. BabelBias measures both within a single corpus --- the
270-of-270 YandexGPT refusal sweep on the silence side, and
the EN-ingroup-pull plus stance-judge results on the content side ---
across a wider provider matrix than either of the comparables.

{\sloppy
\citet{li2024borderlines} introduce \emph{BORDERLINES}, a multilingual
benchmark of 726 multiple-choice questions over 251 disputed territories
in 49 languages, and use a family of \emph{Concurrence Scores} to
measure the discrete
factual flip in a model's claimed controller of a disputed territory
across query languages. They evaluate GPT-4, GPT-3, BLOOM, and BLOOMZ,
with GPT-4 case studies on Crimea, Taiwan, and the Golan Heights. Vanilla
GPT-4 answers ``Russia'' for Crimea in Russian and ``Ukraine'' in
Ukrainian --- a discrete cross-lingual flip in the \emph{opposite
direction} from the EN-ingroup pull reported in
Chapter~\ref{ch:experiments}. The two findings are not
contradictory: cosine in dense embedding space and discrete answer
choice in MCQ are orthogonal measurements, and a model can answer
``Russia'' in Russian while still being closer to the English Wikipedia
anchor in sentence-embedding space. Reconciling the two --- and
integrating the additional Ukrainian-ward stance signal reported in
Chapter~\ref{ch:experiments} --- is a discussion-chapter
contribution.\par}

The provider ecosystem audited in
Chapter~\ref{ch:experiments} contains fifteen providers across seven
regulatory regimes, fourteen of which return analysable prose (the
full matrix is specified in \S\ref{sec:protocol}). The
thesis-relevant claim is not
that these models are similar, but that they share a common epistemic
role: producing language-conditional summaries of contested historical
events for non-expert readers. No prior published work tests
cross-lingual contested-event behaviour across all seven regulatory
regimes simultaneously.

\section{Position of this thesis}\label{sec:novelty}

The literature reviewed above shows that BabelBias is best understood as
a rigorous, multi-regulatory-regime extension of four published findings
--- \citet{durmus2023globalopinionqa}'s Linguistic-Prompting null on
opinions, \citet{buyl2024llm}'s creator-region ideology mapping on
political persons, \citet{urman2025silence}'s cross-lingual guardrail
audit on Western chatbots, and \citet{smirnov2026language}'s
single-document qualitative existence proof of language-conditioned
ideological framing --- rather than a fresh discovery of a previously
unobserved phenomenon. The contribution is the corpus,
the methodology stack, and the methodological controls that the
published partial findings did not have. Five distinct contributions
distinguish the work below.

\begin{enumerate}
  \item \textbf{Multilingual LLM responses scaled and quantified.} The
    measurement target is the LLM's free-form natural-language response
    to a contested-event question, embedded into the same shared
    semantic space as the per-language Wikipedia lead-section anchors.
    This shifts the unit of analysis one level downstream of the prior
    Wikipedia-bias literature \citep{oeberst2020wikipedia,
    samoilenko2017timelines}, and one level upstream of the
    forced-choice MCQ used in \citet{li2024borderlines} or the Likert
    ratings used in \citet{buyl2024llm}. \citet{smirnov2026language}
    establishes the qualitative existence proof on a single contested
    document; this thesis scales the same response-level approach to
    fourteen providers, three languages including English as a control,
    and an imaginary-conflict falsification probe, with the framing
    divergence quantified as anchor-cosine distance rather than
    interpreted via discourse analysis.

  \item \textbf{Provider matrix spanning seven regulatory regimes.} Most
    LLM-bias work tests three to five frontier US-trained models
    \citep{hartmann2023chatgpt, rozado2024,
    durmus2023globalopinionqa}; \citet{buyl2024llm} expands to
    nineteen models across six UN languages but does not include
    Ukrainian as a prompt language and excludes TAIDE and ALLaM.
    \citet{pacheco2025echoes} compares ChatGPT to DeepSeek;
    \citet{urman2025silence} audits three Western chatbots;
    \citet{li2024borderlines} tests four models. BabelBias tests
    fifteen providers from seven regulatory regimes simultaneously
    --- US frontier, Chinese CAC-regulated, Cohere multilingual
    (Canadian), Israeli (Jamba), Saudi state-research (ALLaM),
    Taiwanese state-counter (TAIDE), and Russian state-aligned
    (YandexGPT) --- on the same EN/RU/UK
    contested-event corpus. The provider matrix is the primary
    empirical contribution; no prior published study has the breadth.

  \item \textbf{Methodology-robust by construction.} Findings that
    rest on a single embedder are pre-emptively challenged with the
    four-embedder sweep (\S\ref{sec:exp_015}) including the
    Russian-trained Yandex embedder. The EVoC
    unsupervised-clustering analysis recovers the (question $\times$
    language) cell structure across all four embedders, providing
    convergent validity from a method independent of the anchor
    construction (though sharing the embedding space, so
    embedder-level artefacts are common to both).
    The remaining methodological vulnerability --- the prediction by
    \citet{qi2023crosslingual} that Russian--Ukrainian subword
    vocabulary overlap alone could produce the observed
    UK-clusters-with-RU-anchor finding --- is addressed by the
    tokenizer-overlap control of Chapter~\ref{ch:experiments}.

  \item \textbf{Refusal as primary evidence with embedding-space
    signature.} \citet{buyl2024llm} silently filter ``invalid''
    responses; \citet{durmus2023globalopinionqa} and
    \citet{hartmann2023chatgpt} force option-choice; only
    \citet{urman2025silence} treat refusal as a primary signal, and
    they explicitly leave Russian-creator providers as future work.
    BabelBias closes that gap (the YandexGPT 270/270 refusal sweep
    reported in Chapter~\ref{ch:experiments}) and goes further
    methodologically by characterising
    the refusal embedding-space signature via EVoC clustering,
    recovering the structural inheritance pattern that
    \citet{urman2025silence} predict but cannot verify in their
    Western-only data.

  \item \textbf{Four-way metric reconciliation.} Four complementary
    measurements yield four distinct directional signals on the same
    conflict: continuous anchor-cosine pulls EN-ward; paired-edit
    stance projection compresses toward zero, placing the
    Russo-Ukrainian pair at the floor of the conflicts measured; a
    judged free-form-stance metric on imaginary content tilts
    UK-ward; and \citet{li2024borderlines}'s discrete MCQ
    measurement on the same Crimea item flips RU-ward in vanilla
    GPT-4. Reconciling cosine, stance, judge, and MCQ measurements
    is itself a discussion-chapter contribution that no single prior
    method exposes.
\end{enumerate}

\subsection*{What this thesis does not claim}

This thesis does \emph{not} claim to discover that prompting LLMs in
Russian or Ukrainian fails to elicit Russian or Ukrainian framing ---
that is \citet{durmus2023globalopinionqa}'s Linguistic-Prompting
result, replicated here at the event-narrative level. It does not
claim to discover that LLM ideology depends on the creator's
geopolitical region --- that is \citet{buyl2024llm}'s finding,
replicated here on contested events. It does not claim to discover
that Western chatbots inherit Russian-state censorship templates
cross-lingually --- that is \citet{urman2025silence}'s finding,
extended here to the Russian-creator limit case. It does not claim
to discover that prompt language alone changes the ideological framing
of LLM output on contested political content --- that is
\citet{smirnov2026language}'s qualitative existence proof on a single
document, scaled here across many events and providers, with an
English control and an imaginary-conflict falsification. The empirical
headline numbers in Chapter~\ref{ch:experiments} should be read as
rigorous, controlled, multi-provider extensions of these published
findings, not as fresh discoveries.
